%
%      Brno University of Technology
%      Faculty of Information Technology
%
%      MSc Thesis	2009/2010
%
%      Design of S-Boxes Using Genetic Algorithms
%
%
%      Bedrich Hovorka
%
This short manual is intended for those who want to conduct their own experiments in evolutionary
S-box search.
For proper (parallel) operation, it is necessary to have a compiled modified version of the GAlib library.
For sequential execution, the original version of the library has also been tested.
After installing the library, the dynamic library for S-box search
can be compiled using the \emph{make} command. If for some reason GAlib cannot be installed and can only be compiled
in a directory, the path to this library must be set in the makefile using the -I and -L switches.
(In the submitted form, the source codes are in the state in which they ran on the Merlin server, so the makefile is already configured.)

The objects described below are accessible in (I)Python using the command:
\begin{verbatim}
  from evolution import *
\end{verbatim}
as can be seen in the experiments.py file.
The conducted experiments are also listed as examples in this file.

\subsubsection{Defining the Genome}
The Genome object constructor is used to select the representation, to which the type and
objective function for single-criterion optimization must be passed:
\begin{verbatim}
  Genome(chromosomeType, criterion, [specific])
\end{verbatim}
where \verb|chromosomeType| can be\\
\verb|ChromosomeType.BINARY| for binary representation; \verb|specific| are number of input bits, number of output bits\\
\verb|ChromosomeType.PERMUTATION| for representation as a list of outputs; \verb|specific| permutation requirement, number of input bits,
number of output bits\\
\verb|ChromosomeType.CGP| for representation as an acyclic graph of gates; \verb|specific| number of input bits,
number of output bits, number of columns, number of rows, maximum number of mutations per element\\
\verb|ChromosomeType.SOFTWARE_IMPL| for representation as a sequence of i686 triple-instructions;\\
\verb|specific| number of 4-bit groups at input and output simultaneously\\

and \verb|criterion| \\
(for multi-criteria optimization it is ignored, but \verb|CriterionFunction.NONE_FITNESS| can be specified)
\verb|CriterionFunction.BENT_AND_MOSAC| for bent score criterion,\\
\verb|CriterionFunction.LP_MAX| for linear probability, \\
\verb|CriterionFunction.DP_MAX| for differential probability,\\
\verb|CriterionFunction.BIJECTIVE_SCORE| for bijective score,\\
\verb|CriterionFunction.SAC| for Strict Avalanche Criterion order,\\
\verb|CriterionFunction.BF| for branching factor,\\
\verb|CriterionFunction.POLYNOMIAL_DEGREE| for algebraic polynomial degree

\subsubsection{Defining Search Parameters and Algorithm Selection}
A new search is created by calling the constructor:
\begin{verbatim}
Searching(algorithm, genome, popSize, generations, [specific])
\end{verbatim}
\verb|algorithm| is\\
\verb|"Ga"| for single-criterion genetic algorithm (genetic programming); \\ \verb|specific| mutation probability,
 crossover probability, elitism enabled\\
\verb|"Vega"| for VEGA algorithm; \verb|specific| mutation probability,
 crossover probability, list of criteria\\
\verb|"Spea"| for SPEA algorithm; \verb|specific| mutation probability,
 crossover probability, list of criteria\\
\verb|"Eda"| for EDA algorithm; \verb|specific| type (\verb|False| for UMDA, \verb|True| for BMDA)\\
\verb|"Random"| for random search; population size means number of stored best results\\
\verb|"ParallelRandom"| for parallel random search; \verb|specific| mutation probability\\

\verb|genome| is the representation created according to the rules above
and \verb|popSize| respectively \verb|generations| is the population size respectively number of generations.


\subsubsection{Running Evolution and Obtaining Results}
The Searching class contains:
\begin{itemize}
  \item instance method  \verb|simpleEvolve()| -- runs a single experiment
  \item class method \verb|parallelSearching(list)| -- performs parallel execution of a \verb|list| of Searching objects
  \item instance method \verb|statistics()| -- after evolution completes, returns a TStatistics object
  \item instance method \verb|bestPopulationStrings()| -- returns a list of string representations of all \verb|statistics().nBestGenomes| best individuals throughout the evolution
\end{itemize}

The TStatistics class contains:
\begin{itemize}
  \item \verb|online| -- average score throughout the evolution
  \item \verb|offlineMax| -- average of maximum scores
  \item \verb|offlineMin| -- average of minimum scores
  \item \verb|bestPopulationScores| -- scores of all \verb|nBestGenomes| best individuals throughout the evolution
  \item \verb|bestPopulationCriterionsValues| -- values of individual criteria of all \verb|nBestGenomes| best individuals throughout the evolution
  \item \verb|bestPopulationOutputs| -- outputs of all \verb|nBestGenomes| best individuals throughout the evolution
  \item \verb|maxEver| -- highest score achieved
  \item \verb|minEver| -- lowest score achieved
  \item \verb|generation| -- last generation
  \item \verb|convergence| -- evolution convergence
  \item \verb|selections| -- number of selections performed throughout the evolution
  \item \verb|crossovers| -- number of crossovers performed throughout the evolution
  \item \verb|mutations| -- number of mutations performed throughout the evolution
  \item \verb|replacements| -- number of replacements performed throughout the evolution
  \item \verb|nBestGenomes| -- number of stored best individuals\\
\end{itemize}

Items \verb|bestPopulationScores| etc. are of type \verb|numpy.Array|.

\subsubsection{Performing Symbolic Regression from Previous Search Results}
For symbolic regression, it is necessary to create a specialized genome:
\begin{verbatim}
  SymbolicalRegressionGenome(chromosomeType, list, [specific])
\end{verbatim}
where \verb|list| is the list of required outputs and must match the S-box parameters.
Other parameters are identical to the \verb|Genome| class.
